Injection flaws, cross-site scripting, and insecure data handling remain among the most frequently reported weakness classes in industry surveys \cite{owaspTop10_2021,mitreCWE}. Two families of tools attempt to catch these problems before they reach production. Static Application Security Testing (SAST) tools such as Semgrep and CodeQL scan source code for known patterns and fit naturally into CI/CD pipelines \cite{bessey2010billion,chess2004staticanalysis}. More recently, Large Language Models (LLMs) have been applied to vulnerability detection and repair suggestion \cite{chen2021codex,openai2023gpt4}.

Each family has a gap that the other fills. SAST tools are deterministic and privacy-preserving but shallow: they match syntax without understanding context, produce no repair guidance, and generate enough false positives that developers routinely ignore them \cite{johnson2013don}. LLM-based assistants can reason about context, explain findings, and suggest concrete fixes --- but the dominant offerings require transmitting source code to cloud servers, which is unacceptable for organizations bound by confidentiality requirements or regulations such as GDPR \cite{gdpr2016}. No existing tool provides all three: context-aware detection, actionable repair, and strict source-code privacy.

This thesis addresses that gap. It presents Code Guardian, a VS Code extension that performs LLM-powered vulnerability analysis entirely on the developer's machine. The following sections describe the problem in detail, present a motivating scenario, define the scope and objectives, and outline the rest of the thesis.

\section{Problem Description}

Static Application Security Testing (SAST) tools such as Semgrep and CodeQL have been the industry's primary defense against implementation-level vulnerabilities for over a decade \cite{bessey2010billion,chess2004staticanalysis}. They are deterministic, fast, and integrate well into CI/CD pipelines. But they have two blind spots that limit their practical impact.

First, they are shallow. Rule-based scanners match syntactic patterns but cannot explain \emph{why} a finding matters in context. They cannot tell whether a dangerous-looking sink is already guarded by upstream validation, so they flag it regardless. Developers learn to distrust and eventually ignore the output, especially under delivery pressure \cite{johnson2013don,christakis2016developers}. Second, SAST tools flag problems but almost never suggest how to fix them. The developer must research the vulnerability class, find the right mitigation, and apply it --- a gap that further reduces the chance that findings are acted on.

Large Language Models (LLMs) can fill both gaps: they reason about code context, explain findings, and generate repair suggestions \cite{chen2021codex,openai2023gpt4}. The catch is privacy. Most LLM-powered security tools depend on cloud inference, which means the developer's source code leaves the machine. For organizations with confidentiality constraints, regulatory obligations such as GDPR, or air-gapped environments, this rules out the tool entirely \cite{carlini2021extracting,gdpr2016}.

The result is a three-way trade-off that no current tool resolves. SAST provides determinism and privacy but not contextual reasoning or repair. Cloud LLM assistants provide reasoning and repair but not privacy. A developer who needs all three has no off-the-shelf option. This thesis explores whether a local, modular design --- keeping inference on the developer's machine, grounding model reasoning in curated security knowledge, and integrating into the editor with developer-controlled remediation --- can close that gap.

\section{Motivating Scenario}

To illustrate these challenges in a realistic development setting, consider a typical workflow at a fintech company building a Node.js backend.

A developer is implementing an Express.js endpoint that retrieves user account details. The handler reads a user ID from the request query string and constructs a SQL query to look up the account. The code is straightforward, the tests pass, and the pull request is approved by a colleague:

\begin{lstlisting}[caption={Vulnerable Express.js endpoint with SQL injection}, label={lst:intro-scenario}]
app.get('/api/account', (req, res) => {
  const userId = req.query.id;
  const query = "SELECT * FROM accounts WHERE id = '" + userId + "'";
  db.query(query, (err, results) => {
    res.json(results);
  });
});
\end{lstlisting}

The query is built by string concatenation. An attacker who sends \texttt{id=' OR '1'='1} retrieves every account in the database. No one notices the vulnerability until a penetration test three months later.

This scenario reveals four gaps in existing security tooling:

\begin{itemize}
  \item \textbf{A linter stays silent.} ESLint with security plugins does not flag this pattern because the concatenation is syntactically valid JavaScript. The tool has no concept of tainted input or SQL semantics.
  \item \textbf{A SAST scanner gives a generic warning.} Semgrep might flag string concatenation in a query, but the warning says ``potential SQL injection'' without explaining whether the specific data flow from \texttt{req.query} to \texttt{db.query} is actually dangerous, or how to fix it.
  \item \textbf{A cloud-based LLM assistant explains and repairs --- but requires sending the code offsite.} A tool like GitHub Copilot could identify the injection, explain the tainted data flow, and suggest replacing the concatenation with a parameterized query (\texttt{db.query("SELECT * FROM accounts WHERE id = ?", [userId])}). But this requires transmitting the source code to an external server.
  \item \textbf{No tool provides all three: context-aware detection, actionable repair, and privacy.} The developer needs an assistant that can reason about the vulnerability in context, suggest a concrete fix, and do so without the code leaving the machine.
\end{itemize}

Now imagine a local AI assistant integrated into VS Code. As the developer writes the endpoint, the assistant analyzes the enclosing function, identifies the unsanitized \texttt{req.query.id} flowing into a SQL string, and produces a finding:

\begin{quote}
\textbf{SQL Injection (CWE-89):} User-controlled input \texttt{req.query.id} is concatenated into a SQL query without sanitization. Replace with a parameterized query: \texttt{db.query("SELECT * FROM accounts WHERE id = ?", [userId])}.
\end{quote}

The finding appears as an inline diagnostic in the editor. The developer hovers to read the explanation, clicks the quick-fix action to preview the parameterized query, and applies the patch --- all without the code leaving the machine. This is what Code Guardian provides: contextual detection, grounded explanation, and developer-controlled repair, running entirely on localhost.

\section{Scope of the Thesis}
\label{sec:intro-scope}

This thesis focuses on \textbf{JavaScript/TypeScript} projects in \textbf{Visual Studio Code}. The primary privacy goal is \textbf{no source-code exfiltration}: all code and IDE context are processed exclusively by local services running on the developer's machine. The system is designed as a VS Code extension that connects to a local Ollama instance for LLM inference, with an optional retrieval-augmented generation (RAG) module backed by a local HNSW vector store for security knowledge grounding.

The technical scope includes local LLM inference via Ollama, function-level code scoping, structured JSON output enforcement, and retrieval-augmented prompting with curated security knowledge bases (CWE descriptions, OWASP guidance). No model fine-tuning or training is performed; the focus is on system-level orchestration --- prompt design, output parsing, caching, and IDE integration --- using off-the-shelf local models served through Ollama \cite{ollamaDocs}.

The evaluation is limited to the accuracy, consistency, and latency of the structured output produced from code snippets in a controlled test harness. This thesis does not cover areas such as multi-language support beyond JavaScript/TypeScript, real-time collaborative editing, endpoint hardening, hardware-level attacks, or enterprise identity and network controls. Optional knowledge-base refresh may fetch \emph{public} vulnerability metadata (for example CVE/CWE descriptions) and can be disabled for fully offline operation.

\subsection*{Research Questions}

The thesis addresses three research questions:
\begin{itemize}
  \item \textbf{RQ1 (Feasibility):} To what extent can a locally deployed LLM integrated into VS Code provide useful vulnerability detection and repair suggestions without transmitting source code to external services?
  \item \textbf{RQ2 (Grounding):} How does retrieval-augmented prompting influence detection quality and the qualitative grounding of explanations compared to an LLM-only configuration?
  \item \textbf{RQ3 (Practicality):} What performance mechanisms (scoping, caching, debouncing) are necessary to make LLM-based security feedback usable within interactive IDE workflows?
\end{itemize}

\subsection*{Threat Model and Trust Assumptions}

The main threat addressed is unintended disclosure of proprietary code through external analysis services. Local inference reduces this risk surface but does not remove all security risks.

\begin{table}[H]
  \centering
  \caption{Threat model summary for Code Guardian.}
  \label{tab:intro-threat-model}
  \footnotesize
  \renewcommand{\arraystretch}{1.15}
  \begin{tabularx}{\textwidth}{p{0.19\textwidth}p{0.28\textwidth}X}
    \toprule
    Threat class & Assumption / attack path & Design response in this thesis \\
    \midrule
    Source-code exfiltration & External API calls could expose proprietary code or derived context. & Local inference only; prompts and code remain on localhost in evaluated workflows. \\
    Prompt injection & Attacker-controlled code/comments attempt to override analysis instructions \cite{owaspLLMTop10_2023}. & Strict JSON contract, defensive parsing, and user approval before applying any repair. \\
    Retrieval poisoning & Low-quality or malicious knowledge entries reduce grounding quality. & Curated local knowledge sources and explicit retrieval scope; stronger provenance controls left for future work. \\
    Local host compromise & Malware or untrusted local users access local code, prompts, or caches. & Out of scope; assumes trusted host and OS-level hardening. \\
    \bottomrule
 \end{tabularx}
\end{table}

\section{Objectives}
\label{sec:intro-objectives}

The main objective of this thesis is to design and evaluate an IDE-integrated assistant --- called Code Guardian --- that provides contextual vulnerability detection and repair suggestions while preserving source-code confidentiality. The system aims to overcome the limitations of existing approaches by combining local LLM inference with retrieval-augmented generation in a privacy-preserving architecture that runs entirely on the developer's machine.

To achieve this goal, the following three specific objectives are defined:

\begin{enumerate}
  \item \textbf{Design a privacy-preserving VS Code extension for local vulnerability analysis} that separates the process into a two-stage pipeline: detection followed by developer-controlled repair suggestion. The extension enforces structured JSON output for machine-checkable diagnostics, integrates an optional RAG module backed by curated security knowledge (CWE descriptions, OWASP guidance), and implements performance mechanisms --- function-level scoping, LRU caching, and debounced triggering --- for practical IDE interaction latency.

  \item \textbf{Evaluate system performance across local model configurations} using a documented evaluation harness with 101 test cases spanning 14 CWE categories. Five local models are tested in both LLM-only and LLM+RAG modes, reporting precision, recall, F1, false-positive rate, JSON parse success, and end-to-end latency. Results from all configurations are compared to identify the most effective trade-off between detection quality, output robustness, and response time.

  \item \textbf{Assess the practical impact of retrieval-augmented prompting} by comparing LLM-only and LLM+RAG configurations to determine whether grounding model reasoning in external security knowledge improves detection quality and explanation grounding, or whether the added retrieval context introduces noise that degrades specific model configurations.
\end{enumerate}

These objectives guide the development and validation of a system that is not only technically functional but also practical for interactive developer workflows. Chapter~\ref{chap:evaluation} identifies the strongest configuration among those tested and reports its trade-offs across accuracy, consistency, latency, and privacy.

\section{Thesis Outline}
\label{sec:intro-structure}

Following this introduction, the thesis begins with a review of related work on static analysis tools, LLM-based vulnerability detection, and retrieval-augmented generation for secure coding. The literature is analyzed to identify current limitations and to derive requirements for a system that combines contextual reasoning, repair generation, and source-code privacy. This provides the conceptual foundation for understanding the need for a tool like Code Guardian.

The next chapter introduces the overall system concept. It presents the architecture of Code Guardian together with key design decisions such as the two-stage detection-then-repair pipeline, structured JSON output enforcement, and the optional RAG module. After outlining the concept, the implementation chapter describes the VS Code extension, technology choices, and how core components --- such as the analyzer, RAG manager, analysis cache, and diagnostic rendering --- were developed, with an emphasis on modularity and local execution.

The evaluation chapter examines system performance across five local models in both LLM-only and LLM+RAG configurations, using metrics such as precision, recall, F1, false-positive rate, JSON parse success, and latency. The thesis concludes by summarizing key findings, discussing limitations, and outlining opportunities for future work. Appendices provide supplementary material including prompt templates, detailed evaluation data, implementation details, case studies, and privacy verification evidence.
